\section{Resource Demo and Case Study}

The demo is designed as a diagnostic-resource check rather than a leaderboard.
It asks whether \tool can ingest real public \sid artifacts from both a
canonical residual-quantization line and a recent recommender-native line, then
produce comparable audit tables with explicit caveats. GRID provides the
open RQ-KMeans-style path~\cite{ju2025grid}; ReSID provides the GAOQ line and
processed Musical artifacts~\cite{liang2026resid}. Sanity tokenizers are kept
as controlled references for metric sensitivity.

\begin{table}[t]
\caption{Same-item Musical diagnostic case study. Lower collision is better;
D3 is a co-occurrence neighborhood proxy.}
\label{tab:musical-diagnostic}
\small
\setlength{\tabcolsep}{3.5pt}
\begin{tabular}{@{}lrrrrr@{}}
\toprule
System & Unique & Dup. & Full coll. & D3 L1 & Tail cap. \\
\midrule
GRID feature-text & 3,749 & 0.8421 & 0.9769 & 0.0552 & 0.3695 \\
ReSID GAOQ & 23,742 & 0.0000 & 0.0000 & 0.1535 & 1.0000 \\
\bottomrule
\end{tabular}
\end{table}

Table~\ref{tab:musical-diagnostic} compares two exports on the same 23,742
Musical items. The GRID row is intentionally labeled as feature-text: it uses
processed feature text from the ReSID artifact path and the official
MiniBatchKMeans-style residual quantization module, not a faithful raw-text
TIGER reproduction. Under that controlled setup, the row exposes high collision
pressure and limited tail capacity. The ReSID/GAOQ row exports one full code
per item in this bounded run, yielding zero full-code collisions and higher
D3-L1 collaborative prefix recovery.

This contrast is useful precisely because the table is not a downstream
ranking claim. It shows that the same audit interface can surface different
artifact mechanisms: compact residual-k-means codes can collapse many items
into shared full \sid{}s; recommender-native GAOQ can avoid full collisions in
the exported mapping; and the head-tail table makes the capacity allocation
visible rather than inferred from item popularity. Additional artifact tables
cover sanity controls, GRID scale checks, DACT churn, and a MovieLens schema
smoke, but those are placed in the repository rather than the four-page PDF.
